\section*{Model Formulation}

We use the simple dual-earner model from lecture~6, based on the framework of \cite{borella2023}. The model makes three key simplifications relative to the full model: there are
no savings, the couple cannot divorce, and the environment is deterministic. The model parameters used are given in Table \ref{tab:Model Parameters}

\subsection*{Recursive Formulation}

\begin{align}
    V_t(K_{1,t},K_{2,t})
    &= \max_{h_{1,t},\,h_{2,t}} \;
       U(c_t,\,h_{1,t},\,h_{2,t})
       + \beta\, V_{t+1}(K_{1,t+1},\,K_{2,t+1}) \\[6pt]
    c_t &= \sum_{j=1}^{2} w_{j,t}\,h_{j,t}
           - T(w_{1,t}\,h_{1,t},\; w_{2,t}\,h_{2,t}) \\[6pt]
    \log w_{j,t} &= \alpha_{j,0} + \alpha_{j,1}\,K_{j,t},
        \quad j \in \{1,2\} \\[6pt]
    K_{j,t+1} &= (1-\delta)\,K_{j,t} + h_{j,t},
        \quad j \in \{1,2\}
\end{align}


\subsection*{Preferences}

Preferences are the sum of individual utilities:
\begin{equation}
    U(c_t,\,h_{1,t},\,h_{2,t})
    = 2\,\frac{(c_t/2)^{1+\eta}}{1+\eta}
      - \rho_1\,\frac{h_{1,t}^{1+\gamma}}{1+\gamma}
      - \rho_2\,\frac{h_{2,t}^{1+\gamma}}{1+\gamma}
\end{equation}

\subsection*{Taxes}

Under \textbf{joint taxation} (baseline), taxes are on the household level:
\begin{equation}
    T(Y_1,Y_2) = \bigl(1 - \lambda(Y_1+Y_2)^{-\tau}\bigr) \cdot (Y_1+Y_2)
\end{equation}

The reform of interest is \textbf{individual taxation}, where each member is taxed separately:

\begin{equation}
    T^{\text{indiv}}(Y_1,Y_2)
    = \bigl(1 - \lambda^{\text{indiv}}\,Y_1^{-\tau^{\text{indiv}}}\bigr)\,Y_1
    + \bigl(1 - \lambda^{\text{indiv}}\,Y_2^{-\tau^{\text{indiv}}}\bigr)\,Y_2
\end{equation}
